\documentclass[aspectratio=169,10pt]{beamer}

% Shared Beamer setup for Fall 2026 course slides.
% Clean Cal Poly-inspired theme: no custom class files or uncommon packages.
\mode<presentation>{
  \usetheme{default}
  \usefonttheme{professionalfonts}
}

\definecolor{CalPolyGreen}{HTML}{154734}
\definecolor{CalPolyGold}{HTML}{BD8B13}
\definecolor{CalPolySage}{HTML}{7FA28F}
\definecolor{CalPolyMint}{HTML}{DDF1E7}
\definecolor{CalPolySand}{HTML}{A29F86}
\definecolor{CalPolyBeige}{HTML}{EFEEDE}
\definecolor{DeckInk}{HTML}{1F2933}
\definecolor{DeckMuted}{HTML}{5F6B7A}
\definecolor{DeckSoft}{HTML}{F7F7F5}

\setbeamersize{text margin left=0.55in,text margin right=0.55in}
\setbeamercolor{normal text}{fg=DeckInk,bg=white}
\setbeamercolor{structure}{fg=CalPolyGreen}
\setbeamercolor{title}{fg=white}
\setbeamercolor{frametitle}{fg=CalPolyGreen,bg=white}
\setbeamercolor{block title}{bg=CalPolySage,fg=white}
\setbeamercolor{block body}{bg=CalPolyMint,fg=black}
\setbeamercolor{alerted text}{fg=CalPolyGold}

\setbeamerfont{title}{size=\Huge,series=\bfseries}
\setbeamerfont{frametitle}{size=\Large,series=\bfseries}
\setbeamerfont{block title}{size=\normalsize,series=\bfseries}

\setbeamertemplate{navigation symbols}{}
\setbeamertemplate{itemize item}{\color{CalPolyGold}$\blacktriangleright$}
\setbeamertemplate{itemize subitem}{\color{CalPolyGreen}$\bullet$}
\setbeamertemplate{footline}{
  \hfill{\color{DeckMuted}\scriptsize\insertframenumber/\inserttotalframenumber}
  \hspace{0.18in}\vspace{0.12in}
}
\setbeamertemplate{frametitle}{
  \vspace{0.18in}
  \hspace{0.02in}\insertframetitle\par
  \vspace{0.08in}
  \noindent\makebox[\linewidth][l]{%
    {\color{CalPolyGold}\rule{0.55in}{2pt}}%
    {\color{CalPolyGreen}\rule{\dimexpr\linewidth-0.55in\relax}{2pt}}%
  }
}

\newcommand{\CourseNumber}{Course}
\newcommand{\CourseTitle}{Course Title}
\newcommand{\LectureNumber}{0}
\newcommand{\LectureTitle}{Lecture Title}
\newcommand{\LectureDate}{Date}
\newcommand{\CourseUnit}{Unit}

\newcommand{\coursenumber}[1]{\renewcommand{\CourseNumber}{#1}}
\newcommand{\coursetitle}[1]{\renewcommand{\CourseTitle}{#1}}
\newcommand{\lecturenumber}[1]{\renewcommand{\LectureNumber}{#1}}
\newcommand{\lecturetitle}[1]{\renewcommand{\LectureTitle}{#1}}
\newcommand{\lecturedate}[1]{\renewcommand{\LectureDate}{#1}}
\newcommand{\courseunit}[1]{\renewcommand{\CourseUnit}{#1}}

\author{Kirk Duran}
\institute{California Polytechnic State University, San Luis Obispo}
\date{\LectureDate}

\newcommand{\makedecktitle}{
  \begingroup
  \setbeamercolor{background canvas}{bg=CalPolyGreen}
  \begin{frame}[plain]
    \vfill
    \begin{center}
      {\usebeamerfont{title}\color{white}\LectureTitle\par}
      \vspace{0.18in}
      {\Large\color{white}Lecture \LectureNumber\par}
    \end{center}
    \vfill
    \noindent{\color{white}\rule{\linewidth}{1.2pt}}\par
    \vspace{0.08in}
    \noindent\colorbox{CalPolyGold}{\makebox[0.24\linewidth][c]{\color{white}\Large\CourseNumber}}
    \hspace{0.12in}{\color{white}\Large Kirk Duran}
  \end{frame}
  \endgroup
}

\newenvironment{keyidea}{\begin{block}{Key idea}}{\end{block}}
\newenvironment{activity}{\begin{block}{In class}}{\end{block}}
\newenvironment{cpdefinition}{\begin{block}{Definition}}{\end{block}}
\newenvironment{exampleblockcp}{
  \begingroup
  \setbeamercolor{block title}{bg=CalPolySand,fg=white}
  \setbeamercolor{block body}{bg=CalPolyBeige,fg=black}
  \begin{block}{Example}
}{
  \end{block}
  \endgroup
}

\newcommand{\imageplaceholder}[2][0.3in]{%
  \noindent\makebox[\linewidth][c]{%
    \fcolorbox{CalPolySage}{CalPolyBeige}{%
      \parbox[c][#1][c]{0.86\linewidth}{%
        \centering
        {\color{DeckMuted}\scriptsize Image placeholder: }%
        {\color{DeckInk}\scriptsize\ttfamily\detokenize{#2}}%
      }%
    }%
  }%
}

\newcommand{\sectionframe}[1]{
  \begingroup
  \setbeamercolor{background canvas}{bg=CalPolyGreen}
  \begin{frame}[plain]
    \vfill
    {\color{white}\LARGE\bfseries #1\par}
    \vspace{0.18in}
    {\color{CalPolyGold}\rule{0.22\paperwidth}{3pt}}
    {\color{white}\rule{0.62\paperwidth}{3pt}}
    \vfill
  \end{frame}
  \endgroup
}

\newcommand{\placeholdernote}{
  \begin{block}{Placeholder}
    Replace these bullets with the final lecture material, examples, and in-class activities.
  \end{block}
}


\pdfstringdefDisableCommands{\def\translate#1{#1}}

\graphicspath{{../assets/lecture-17/}}

% If an image file is not yet available, the slide displays a labeled
% placeholder using the intended filename. Place the matching file in
% ../assets/lecture-17/ to replace the placeholder automatically.
\newcommand{\lectureimage}[2][1.75in]{%
  \IfFileExists{../assets/lecture-17/#2}{%
    \includegraphics[width=\linewidth,height=#1,keepaspectratio]{#2}%
  }{%
    \fbox{%
      \parbox[c][#1][c]{0.94\linewidth}{%
        \centering
        \small\textbf{Image Placeholder}\\[0.08in]
        \scriptsize\texttt{\detokenize{#2}}
      }%
    }%
  }%
}

\setbeamersize{text margin left=0.42in,text margin right=0.42in}

\coursenumber{CS 1001}
\coursetitle{Introduction to Computing}
\lecturenumber{17}
\lecturetitle{Lists I --- Sequential Data}
\lecturedate{Fall 2026}
\courseunit{7. Sequential Data}

\begin{document}

\makedecktitle

\begin{frame}{Today}
  \begin{itemize}
    \item Motivate lists as one structure for representing related data.
    \item Define lists as ordered, indexed, mutable sequences of elements.
    \item Construct lists with list literals, including the empty list and mixed-type lists.
    \item Access elements using zero-based positive and negative indexes.
    \item Test membership and locate an element with \texttt{.index()}.
    \item Mutate one element by assigning through an index.
    \item Build a useful memory model, then refine it to match how Python lists are actually represented.
  \end{itemize}

  \begin{keyidea}
    A list gives one name to an ordered sequence of element positions. Each position can be read or updated by its index.
  \end{keyidea}
\end{frame}

\sectionframe{Part 1: Why Sequential Data?}

\begin{frame}{Related Data Is Awkward When Every Value Has Its Own Name}
  \begin{columns}[T,onlytextwidth]
    \begin{column}{0.50\textwidth}
      \begin{block}{Separate variables}
        \texttt{test1 = 90}\\
        \texttt{test2 = 88}\\
        \texttt{test3 = 76}\\
        \texttt{test4 = 92}
      \end{block}

      \begin{itemize}
        \item The values have the same conceptual role: test scores.
        \item Their relationship exists only in the programmer's naming convention.
        \item Adding more related values requires inventing more names.
      \end{itemize}
    \end{column}
    \begin{column}{0.44\textwidth}
      % Intended visual: many separate score variables contrasted with one grouped structure.
      \lectureimage[2.45in]{related-data-variables-vs-list.png}
    \end{column}
  \end{columns}
\end{frame}

\begin{frame}{A List Groups Related Values Under One Name}
  \begin{block}{One sequence of scores}
    \centering
    \Large\texttt{scores = [90, 88, 76, 92]}
  \end{block}

  \begin{itemize}
    \item \texttt{scores} is one variable name.
    \item The value bound to \texttt{scores} is a list containing four elements.
    \item Each element occupies a specific position in the sequence.
  \end{itemize}

  \begin{keyidea}
    The list represents the relationship explicitly: these values belong to one ordered sequence.
  \end{keyidea}
\end{frame}

\begin{frame}{Definition: Python List}
  \begin{cpdefinition}
    A \textbf{list} is an ordered, indexed, mutable sequence of zero or more elements.
  \end{cpdefinition}

  \begin{center}
    \renewcommand{\arraystretch}{1.35}
    \begin{tabular}{l|l}
      \textbf{Property} & \textbf{Meaning} \\\hline
      ordered & element position is part of the list's structure \\
      indexed & a position can be addressed by an integer index \\
      mutable & an existing element can be replaced after creation \\
      sequence & multiple elements are represented in a defined order
    \end{tabular}
  \end{center}
\end{frame}

\begin{frame}{Order Is Part of the List Value}
  \begin{columns}[T,onlytextwidth]
    \begin{column}{0.46\textwidth}
      \begin{block}{List A}
        \centering
        \texttt{[10, 20, 30]}
      \end{block}
    \end{column}
    \begin{column}{0.46\textwidth}
      \begin{block}{List B}
        \centering
        \texttt{[30, 20, 10]}
      \end{block}
    \end{column}
  \end{columns}

  \begin{itemize}
    \item The lists contain the same three integer values.
    \item They do not represent the same sequence because the positions differ.
    \item Consequently, \texttt{[10, 20, 30] == [30, 20, 10]} evaluates to \texttt{False}.
  \end{itemize}

  \begin{keyidea}
    In sequential data, position carries information.
  \end{keyidea}
\end{frame}

\sectionframe{Part 2: Creating Lists}

\begin{frame}{List Literals Use Square Brackets}
  \begin{block}{General form}
    \centering
    \Large\texttt{[element0, element1, element2, ...]}
  \end{block}

  \begin{itemize}
    \item Square brackets delimit the list literal.
    \item Commas separate element expressions.
    \item Each element expression is evaluated and becomes one element of the new list.
  \end{itemize}

  \begin{block}{Examples}
    \texttt{scores = [90, 88, 76]}\\[0.06in]
    \texttt{names = ["Ada", "Grace", "Alan"]}\\[0.06in]
    \texttt{flags = [True, False, True]}\\[0.06in]
    \texttt{empty = []} \quad zero elements and no valid element index
  \end{block}
\end{frame}

\begin{frame}{Python Lists May Contain Mixed Types}
  \begin{block}{Heterogeneous list}
    \centering
    \Large\texttt{record = [42, "hello", True, 3.5]}
  \end{block}

  \begin{itemize}
    \item The elements do not need to share one Python type.
    \item Position remains well-defined even when the element types differ.
    \item Mixed types are legal; whether they are a clear design choice depends on what the positions represent.
  \end{itemize}

  \begin{keyidea}
    Python constrains the structure of the sequence, not the type of every value stored in that sequence.
  \end{keyidea}
\end{frame}

\begin{frame}{Our Working Model: Adjacent Numbered Slots}
  \begin{columns}[T,onlytextwidth]
    \begin{column}{0.42\textwidth}
      % Intended visual: four adjacent boxes containing 90, 88, 76, 92 with indexes 0--3 beneath them.
      \lectureimage[1.80in]{list-adjacent-slots-model.png}
    \end{column}
    \begin{column}{0.54\textwidth}
      \begin{itemize}
        \item Think of a list as a row of adjacent element slots.
        \item Each slot has a unique integer position called its index.
        \item This model makes access and mutation easy to reason about.
      \end{itemize}

      \begin{alertblock}{Modeling note}
        This is our conceptual model for tracing. We will refine the physical memory representation at the end of the lecture.
      \end{alertblock}
    \end{column}
  \end{columns}
\end{frame}

\sectionframe{Part 3: Indexing}

\begin{frame}{An Index Identifies One Position in a Sequence}
  \begin{cpdefinition}
    An \textbf{index} is an integer used to identify a particular element position in a sequence.
  \end{cpdefinition}

  \begin{block}{Example list}
    \centering
    \texttt{fruits = ["apple", "banana", "cherry", "date"]}
  \end{block}

  \begin{center}
    \renewcommand{\arraystretch}{1.3}
    \begin{tabular}{c|c|c|c|c}
      element & apple & banana & cherry & date \\\hline
      index & 0 & 1 & 2 & 3
    \end{tabular}
  \end{center}
\end{frame}

\begin{frame}{Python Uses Zero-Based Indexing}
  \begin{itemize}
    \item The first element is at index \texttt{0}, not index \texttt{1}.
    \item The second element is at index \texttt{1}.
    \item For a list containing \(n\) elements, the last positive index is \(n-1\).
  \end{itemize}

  \begin{block}{Example}
    \texttt{fruits = ["apple", "banana", "cherry", "date"]}\\[0.10in]
    \texttt{fruits[0]} $\longrightarrow$ \texttt{"apple"}\\
    \texttt{fruits[1]} $\longrightarrow$ \texttt{"banana"}\\
    \texttt{fruits[3]} $\longrightarrow$ \texttt{"date"}
  \end{block}
\end{frame}

\begin{frame}{Indexing Is an Expression}
  \begin{block}{Expression}
    \centering
    \Large\texttt{fruits[2]}
  \end{block}

  \begin{enumerate}
    \item Evaluate the name \texttt{fruits} to obtain the list.
    \item Evaluate the index expression \texttt{2}.
    \item Select the element at position \texttt{2}.
    \item The complete indexing expression evaluates to that element's value.
  \end{enumerate}

  \begin{keyidea}
    Just like arithmetic or a function call, an indexing expression produces a value and can participate in a larger expression, such as \texttt{scores[1] + 5}.
  \end{keyidea}
\end{frame}

\begin{frame}{Trace: Reading an Element Does Not Change the List}
  \begin{block}{Code}
    \texttt{scores = [90, 88, 76]}\\
    \texttt{middle = scores[1]}\\
    \texttt{result = middle + 2}
  \end{block}

  \begin{center}
    \renewcommand{\arraystretch}{1.3}
    \begin{tabular}{l|l|l}
      \textbf{Statement} & \textbf{Expression result} & \textbf{Bindings afterward} \\\hline
      \texttt{scores = [90,88,76]} & list value & \texttt{scores} $\rightarrow$ \texttt{[90,88,76]} \\
      \texttt{middle = scores[1]} & \texttt{88} & \texttt{middle} $\rightarrow$ \texttt{88} \\
      \texttt{result = middle + 2} & \texttt{90} & \texttt{result} $\rightarrow$ \texttt{90}
    \end{tabular}
  \end{center}

  \begin{keyidea}
    Accessing an element retrieves a value; it does not mutate the list.
  \end{keyidea}
\end{frame}

\begin{frame}{An Index Must Refer to an Existing Position}
  \begin{block}{List}
    \centering
    \texttt{scores = [90, 88, 76]}
  \end{block}

  \begin{columns}[T,onlytextwidth]
    \begin{column}{0.46\textwidth}
      \begin{block}{Valid}
        \texttt{scores[0]}\\
        \texttt{scores[1]}\\
        \texttt{scores[2]}
      \end{block}
    \end{column}
    \begin{column}{0.46\textwidth}
      \begin{alertblock}{Invalid}
        \texttt{scores[3]}\\[0.05in]
        raises \texttt{IndexError}
      \end{alertblock}
    \end{column}
  \end{columns}

  \begin{keyidea}
    The number of elements is three, but the valid positive indexes are \texttt{0}, \texttt{1}, and \texttt{2}.
  \end{keyidea}
\end{frame}

\begin{frame}{Negative Indexes Count Backward from the End}
  \begin{block}{List}
    \centering
    \texttt{fruits = ["apple", "banana", "cherry", "date"]}
  \end{block}

  \begin{columns}[T,onlytextwidth]
    \begin{column}{0.58\textwidth}
      % Intended visual: same list with positive indexes 0--3 above and negative indexes -4---1 below.
      \lectureimage[1.90in]{positive-negative-indexing.png}
    \end{column}
    \begin{column}{0.38\textwidth}
      \begin{block}{Examples}
        \texttt{fruits[-1]} $\longrightarrow$ \texttt{"date"}\par
        \texttt{fruits[-2]} $\longrightarrow$ \texttt{"cherry"}
      \end{block}
    \end{column}
  \end{columns}
\end{frame}

\begin{frame}{Positive and Negative Indexes Can Name the Same Position}
  \begin{center}
    \renewcommand{\arraystretch}{1.35}
    \begin{tabular}{c|c|c}
      \textbf{Element} & \textbf{Positive index} & \textbf{Negative index} \\\hline
      apple & 0 & -4 \\
      banana & 1 & -3 \\
      cherry & 2 & -2 \\
      date & 3 & -1
    \end{tabular}
  \end{center}

  \begin{keyidea}
    For a list containing \(n\) elements, a valid negative index \(i\) identifies the same position as \(n+i\).
  \end{keyidea}

  \begin{exampleblockcp}
    With four elements, index \texttt{-2} corresponds to \(4 + (-2) = 2\). The valid negative indexes are \texttt{-1} through \texttt{-4}; \texttt{-5} raises \texttt{IndexError}.
  \end{exampleblockcp}
\end{frame}

\begin{frame}{Prediction: Resolve Each Index Before Running Python}
  \begin{block}{Given}
    \centering
    \texttt{values = [12, 24, 36, 48, 60]}
  \end{block}

  \begin{activity}
    Determine the value or error produced by each expression:
    \begin{enumerate}
      \item \texttt{values[0]}
      \item \texttt{values[3]}
      \item \texttt{values[-1]}
      \item \texttt{values[-5]}
      \item \texttt{values[5]}
    \end{enumerate}
  \end{activity}
\end{frame}

\sectionframe{Part 4: Membership and Finding Positions}

\begin{frame}{Membership Asks Whether a Value Is Present}
  \begin{cpdefinition}
    A \textbf{membership test} asks whether a value occurs as an element of a sequence and produces a Boolean result.
  \end{cpdefinition}

  \begin{block}{Examples}
    \texttt{fruits = ["apple", "banana", "cherry"]}\\[0.08in]
    \texttt{"banana" in fruits} $\longrightarrow$ \texttt{True}\\
    \texttt{"orange" in fruits} $\longrightarrow$ \texttt{False}\\
    \texttt{"orange" not in fruits} $\longrightarrow$ \texttt{True}
  \end{block}

  \begin{keyidea}
    Membership answers an existence question; it does not produce the element's position.
  \end{keyidea}
\end{frame}

\begin{frame}{Trace: Membership Produces a Boolean Value}
  \begin{block}{Code}
    \texttt{colors = ["red", "green", "blue"]}\\
    \texttt{target = "green"}\\
    \texttt{found = target in colors}
  \end{block}

  \begin{center}
    \renewcommand{\arraystretch}{1.3}
    \begin{tabular}{l|l}
      \textbf{Expression} & \textbf{Value} \\\hline
      \texttt{target} & \texttt{"green"} \\
      \texttt{colors} & \texttt{["red","green","blue"]} \\
      \texttt{target in colors} & \texttt{True}
    \end{tabular}
  \end{center}

  \begin{keyidea}
    The final binding is \texttt{found} $\rightarrow$ \texttt{True}; the list itself is unchanged.
  \end{keyidea}
\end{frame}

\begin{frame}{\texttt{.index()} Asks Where a Value Occurs}
  \begin{block}{Example}
    \texttt{colors = ["red", "green", "blue"]}\\[0.08in]
    \texttt{colors.index("green")} $\longrightarrow$ \texttt{1}
  \end{block}

  \begin{cpdefinition}
    The list method \texttt{.index(value)} returns the index of the first element equal to \texttt{value}.
  \end{cpdefinition}

  \begin{itemize}
    \item If a value occurs more than once, \texttt{.index()} returns the first matching position.
    \item Example: \texttt{[90, 75, 88, 75].index(75)} $\longrightarrow$ \texttt{1}.
  \end{itemize}

  \begin{keyidea}
    Indexing uses a known position to obtain a value; \texttt{.index()} uses a known value to obtain a position.
  \end{keyidea}
\end{frame}

\begin{frame}{\texttt{.index()} Raises an Error When the Value Is Absent}
  \begin{block}{Example}
    \texttt{colors = ["red", "green", "blue"]}\\[0.08in]
    \texttt{colors.index("orange")}
  \end{block}

  \begin{alertblock}{Result}
    Python raises \texttt{ValueError} because no matching element exists in the list.
  \end{alertblock}

  \begin{keyidea}
    Use a membership test when the immediate question is whether the value exists. Use \texttt{.index()} when the required result is its position.
  \end{keyidea}
\end{frame}

\begin{frame}{Membership and \texttt{.index()} Answer Different Questions}
  \begin{center}
    % Intended visual: split diagram, value -> membership -> bool and value -> .index() -> integer position.
    \lectureimage[1.65in]{membership-versus-index.png}
  \end{center}

  \begin{center}
    \renewcommand{\arraystretch}{1.35}
    \begin{tabular}{l|l|l}
      \textbf{Question} & \textbf{Operation} & \textbf{Result type} \\\hline
      Is it present? & \texttt{value in items} & \texttt{bool} \\
      Where is it? & \texttt{items.index(value)} & \texttt{int}
    \end{tabular}
  \end{center}
\end{frame}

\sectionframe{Part 5: Mutating Individual Elements}

\begin{frame}{Lists Are Mutable}
  \begin{cpdefinition}
    A value is \textbf{mutable} when its state can be changed after the value has been created.
  \end{cpdefinition}

  \begin{block}{Before}
    \centering
    \texttt{scores = [90, 88, 76]}
  \end{block}

  \begin{block}{Indexed assignment}
    \centering
    \Large\texttt{scores[1] = 91}
  \end{block}

  \begin{block}{After}
    \centering
    \texttt{scores} $\longrightarrow$ \texttt{[90, 91, 76]}
  \end{block}
\end{frame}

\begin{frame}{Indexed Assignment Replaces One Existing Element}
  \begin{columns}[T,onlytextwidth]
    \begin{column}{0.49\textwidth}
      \begin{enumerate}
        \item Evaluate the list expression.
        \item Evaluate the index expression.
        \item Evaluate the right-hand-side expression.
        \item Replace the element at the selected position with the resulting value.
      \end{enumerate}
    \end{column}
    \begin{column}{0.47\textwidth}
      % Intended visual: [90,88,76] with index 1 highlighted and changed from 88 to 91.
      \lectureimage[1.65in]{indexed-mutation.png}
    \end{column}
  \end{columns}

  \begin{keyidea}
    The list remains the same list value conceptually, but its element state has changed. A negative index can identify the target as well: \texttt{values[-1] = 99}.
  \end{keyidea}
\end{frame}

\begin{frame}{Reading and Writing Use Similar Syntax but Different Semantics}
  \begin{columns}[T,onlytextwidth]
    \begin{column}{0.46\textwidth}
      \begin{block}{Read}
        \texttt{x = scores[1]}\\[0.10in]
        \texttt{scores[1]} is an expression.\\
        It produces the current element value.
      \end{block}
    \end{column}
    \begin{column}{0.46\textwidth}
      \begin{block}{Write}
        \texttt{scores[1] = 91}\\[0.10in]
        \texttt{scores[1]} is an assignment target.\\
        It identifies the position to replace.
      \end{block}
    \end{column}
  \end{columns}

  \begin{keyidea}
    Context determines whether bracket notation retrieves an element or identifies where assignment should store a new value.
  \end{keyidea}
\end{frame}

\begin{frame}{Invalid Mutation Targets Raise \texttt{IndexError}}
  \begin{block}{Example}
    \texttt{items = [10, 20, 30]}\\[0.08in]
    \texttt{items[3] = 40}
  \end{block}

  \begin{alertblock}{Why this fails}
    Index \texttt{3} does not identify an existing element position. Indexed assignment replaces one existing element; it neither adds a new position nor changes the number of positions in the list.
  \end{alertblock}
\end{frame}

\begin{frame}{Trace: List State Can Change Across Statements}
  \begin{block}{Code}
    \texttt{values = [5, 10, 15]}\\
    \texttt{x = values[1]}\\
    \texttt{values[1] = x + 2}\\
    \texttt{last = values[-1]}
  \end{block}

  \begin{center}
    \renewcommand{\arraystretch}{1.25}
    \begin{tabular}{l|l|l}
      \textbf{After statement} & \textbf{Relevant scalar binding} & \textbf{List state} \\\hline
      list creation & --- & \texttt{[5,10,15]} \\
      \texttt{x = values[1]} & \texttt{x = 10} & \texttt{[5,10,15]} \\
      indexed assignment & \texttt{x = 10} & \texttt{[5,12,15]} \\
      \texttt{last = values[-1]} & \texttt{last = 15} & \texttt{[5,12,15]}
    \end{tabular}
  \end{center}
\end{frame}

\begin{frame}{Prediction: Track Both Position and State}
  \begin{block}{Code}
    \texttt{data = [8, 6, 7, 5]}\\
    \texttt{data[0] = data[-1]}\\
    \texttt{middle = data[2]}\\
    \texttt{data[-2] = 9}
  \end{block}

  \begin{activity}
    Without running Python, determine:
    \begin{enumerate}
      \item the list after the first mutation;
      \item the value bound to \texttt{middle};
      \item the final state of \texttt{data}.
    \end{enumerate}
  \end{activity}
\end{frame}

\sectionframe{Part 6: From the Conceptual Model to Physical Memory}

\begin{frame}{Why the Adjacent-Slot Model Is Useful}
  \begin{columns}[T,onlytextwidth]
    \begin{column}{0.48\textwidth}
      % Intended visual: base address plus fixed-width slot offset leading directly to an indexed slot.
      \lectureimage[2.15in]{indexed-slot-address-model.png}
    \end{column}
    \begin{column}{0.48\textwidth}
      \begin{itemize}
        \item A numbered row of fixed-width slots provides a direct way to locate a position.
        \item Conceptually, position \(i\) is found from the beginning of the row plus an offset for \(i\) slots.
        \item This explains why an index identifies one element position rather than searching by a variable name.
      \end{itemize}
    \end{column}
  \end{columns}
\end{frame}

\begin{frame}{Many Low-Level Arrays Store Homogeneous Values}
  \begin{itemize}
    \item In many programming languages, a conventional array stores elements of one declared type.
    \item A homogeneous element type gives every array slot a predictable representation and size.
    \item Adjacent fixed-size elements therefore support direct address calculation from an index.
  \end{itemize}

  \begin{block}{Conceptual homogeneous array}
    \centering
    \texttt{[10][20][30][40]}\qquad all slots store the same kind of value
  \end{block}

  \begin{keyidea}
    Python preserves the indexed array idea but changes what each physical slot stores.
  \end{keyidea}
\end{frame}

\begin{frame}{A Python List Stores Adjacent References, Not Adjacent Objects}
  \begin{columns}[T,onlytextwidth]
    \begin{column}{0.52\textwidth}
      \begin{itemize}
        \item In CPython, the list maintains a dynamically allocated contiguous array of object references.
        \item Each physical slot stores a reference to a Python object.
        \item The referenced objects themselves may reside elsewhere in memory.
      \end{itemize}

      \begin{keyidea}
        What is physically adjacent is the array of references.
      \end{keyidea}
    \end{column}
    \begin{column}{0.42\textwidth}
      % Intended visual: contiguous reference slots with arrows to separate int, str, bool, float objects.
      \lectureimage[2.65in]{python-list-references.png}
    \end{column}
  \end{columns}
\end{frame}

\begin{frame}{References Explain Why Mixed-Type Lists Work}
  \begin{block}{Python list}
    \centering
    \texttt{record = [42, "hello", True, 3.5]}
  \end{block}

  \begin{itemize}
    \item The integer, string, Boolean, and floating-point objects have different representations.
    \item The list slots do not need to store those object representations directly.
    \item Each slot stores the same kind of thing: a reference to a Python object.
  \end{itemize}

  \begin{keyidea}
    Uniform reference slots allow one Python list to refer to objects of different types.
  \end{keyidea}
\end{frame}

\begin{frame}{Keep the Simpler Model for Ordinary Tracing}
  \begin{columns}[T,onlytextwidth]
    \begin{column}{0.46\textwidth}
      \begin{block}{Tracing model}
        \texttt{[90][88][76]}\\[0.10in]
        Think of indexed values in adjacent numbered positions.
      \end{block}
    \end{column}
    \begin{column}{0.46\textwidth}
      \begin{block}{Implementation model}
        adjacent reference slots\\[0.10in]
        $\downarrow\quad\downarrow\quad\downarrow$\\[-0.02in]
        Python objects elsewhere in memory
      \end{block}
    \end{column}
  \end{columns}

  \begin{keyidea}
    Use the simplest model that correctly predicts the behavior being studied. Refine the model when implementation details become relevant.
  \end{keyidea}
\end{frame}

\begin{frame}{Lecture 17 Mental Model}
  \begin{center}
    \renewcommand{\arraystretch}{1.45}
    \begin{tabular}{l|l}
      \textbf{Concept} & \textbf{Operational question} \\\hline
      list & What ordered sequence of elements is represented? \\
      index & Which position is being addressed? \\
      indexing expression & What value occupies that position? \\
      membership & Does this value occur in the sequence? \\
      \texttt{.index()} & What is the first matching position? \\
      indexed assignment & What existing position is being replaced?
    \end{tabular}
  \end{center}
\end{frame}

\begin{frame}{Takeaways}
  \begin{itemize}
    \item A list is an ordered, indexed, mutable sequence.
    \item List literals use square brackets; \texttt{[]} represents an empty list.
    \item Positive indexes begin at \texttt{0}; negative indexes count backward from the end.
    \item Invalid element positions raise \texttt{IndexError}.
    \item \texttt{in} and \texttt{not in} answer membership questions with Boolean values.
    \item \texttt{.index(value)} returns the first matching position and raises \texttt{ValueError} when no match exists.
    \item Indexed assignment mutates one existing element while preserving the sequence's positions.
    \item For tracing, think in adjacent numbered positions; physically, CPython stores an adjacent array of references to objects.
  \end{itemize}
\end{frame}

\end{document}
